\subsection{Otros Costo de Implementación de Mejora}
\subsubsection{El Costo de Servicio de consumo de Luz}
Se presentara el consumo de energía total que utilizaron los equipos de cómputo durante el tiempo de desarrollo.

\par
\noindent
\begin{minipage}{\linewidth}
    \tituloTablaAPA{tab:costo_energia_laptops}{Consumo de energía eléctrica y costo total por equipos de cómputo}

    \begingroup
    \renewcommand{\baselinestretch}{1}
    \footnotesize
    \setlength{\tabcolsep}{3pt}
    \renewcommand{\arraystretch}{1.25}

    \begin{tabularx}{\linewidth}{@{}p{3.0cm}c c c r r r@{}}
        \toprule
        \textbf{Equipo} & \textbf{Potencia} & \textbf{Horas Uso} & \textbf{Consumo} & \textbf{Pago/Mes (S/)} & \textbf{Total (S/)} & \textbf{Dólares (USD)} \\ \midrule
        Laptop 1 (HP Omen 16) & 65 W & 1,170 h & 76.05 kWh & S/ 7.06 & 49.43 & 14.37 \\ 
        Laptop 2 (ASUS Vivobook 16) & 65 W & 1,170 h & 76.05 kWh & S/ 7.06 & 49.43 & 14.37 \\ 
        Laptop 3 (ASUS Vivobook 14) & 65 W & 1,170 h & 76.05 kWh & S/ 7.06 & 49.43 & 14.37 \\ 
        Laptop 4 (Lenovo ThinkPad X1) & 65 W & 1,170 h & 76.05 kWh & S/ 7.06 & 49.43 & 14.37 \\ 
        Laptop 5 (Lenovo IdeaPad Slim 3) & 65 W & 402 h & 26.13 kWh & S/ 5.66 & 16.98 & 4.94 \\ 
        Laptop 6 (ASUS TUF Gaming F16) & 65 W & 402 h & 26.13 kWh & S/ 5.66 & 16.98 & 4.94 \\ \midrule
        \textbf{Total} & \textbf{390 W} & \textbf{5,484 h} & \textbf{356.46 kWh} & \textbf{--} & \textbf{231.68} & \textbf{67.35} \\ \bottomrule
    \end{tabularx}
    \endgroup

    \notaTabla{Cálculo basado en una potencia nominal promedio de trabajo de 65 W por laptop y una tarifa de energía eléctrica de S/ 0.65 por kWh. Se contempla una jornada laboral de 6 horas diarias durante 195 días (7 meses) para las laptops 1 a 4 y 67 días (3 meses) para las laptops 5 y 6. Tipo de cambio referencial: 1 USD = S/ 3.44 (fecha: 09/11/2026). Elaboración propia.}
\end{minipage}
\finBloqueAPA

\subsubsection{Lista de gastos adicionales}
En este apartado se detallara los gastos adicionales a servicios idenpendientes a la empresa, como energia, internet, entre otros.

\par
\noindent
\begin{minipage}{\linewidth}
    \tituloTablaAPA{tab:gastos_adicionales}{Gastos adicionales en energía eléctrica e internet}

    \begingroup
    \renewcommand{\baselinestretch}{1}
    \footnotesize
    \setlength{\tabcolsep}{4pt}
    \renewcommand{\arraystretch}{1.25}

    \begin{tabularx}{\linewidth}{@{}p{4.5cm}c p{3.0cm}r r@{}}
        \toprule
        \textbf{Descripción} & \textbf{Cantidad} & \textbf{Precio Unitario} & \textbf{Total (S/)} & \textbf{Dólares (USD)} \\ \midrule
        Servicio de Luz (Laptops 1 - 4) & 7 meses & S/ 7.06 / mes (x4) & S/ 197.68 & \$57.47 \\ 
        Servicio de Luz (Laptops 5 - 6) & 3 meses & S/ 7.06 / mes (x2) & S/ 42.36 & \$12.31 \\ 
        Servicio de Internet & 7 meses & S/ 89.90 / mes & S/ 629.30 & \$182.94 \\ \midrule
        \textbf{Total} & & & \textbf{S/ 869.34} & \textbf{\$252.72} \\ \bottomrule
    \end{tabularx}
    \endgroup

    \notaTabla{Cálculo del gasto de luz basado en un costo mensual por laptop de S/ 7.06 (6 horas diarias a S/ 0.65/kWh). El servicio de internet contempla 7 meses de desarrollo del proyecto. Tipo de cambio referencial: 1 USD = S/ 3.44 (fecha: 09/11/2026). Elaboración propia.}
\end{minipage}
\finBloqueAPA